\documentclass[11pt]{article}
\usepackage{fontspec,amsmath,amssymb,amsthm,hyperref}
\usepackage[margin=1in]{geometry}
\setlength{\emergencystretch}{3em}
\title{Iterating concrete stochastic segments to a Wasserstein bound}
\author{LangevinLean mathematical exposition}
\date{}
\begin{document}\maketitle
\noindent Source module: \texttt{KineticLangevin/LpConvergence.lean}. Compile this proof document with XeLaTeX or LuaLaTeX. The accompanying HTML pairs the mathematical argument with the exact Lean statements and records the scope of the formalization.


The numerical phase knots, auxiliary exact segments, reference segments, common raw inputs and convolution noises are explicit. The local coupling inequality is proved internally from these paths. The remaining law inputs state the actual reference Gaussian velocity and stationary position marginals. They do not assume a convergence conclusion.


\section{Accumulate actual update errors and project to laws}


For each step, the auxiliary path begins at the numerical knot, while the reference segment joins consecutive reference knots. The numerical knot update is exactly \texttt{exponentialStep} with innovation \((B(h),\int_0^hB(s)\,ds)\). Rewriting these endpoint identities in the proved L\ensuremath{{}^{2}} one-step theorem produces the recurrence \(e_{k+1}\le e^{-mh/\gamma}e_k+3Mh^2(\sqrt p+e_k)/4\). The source step restriction \(4\gamma Mh\le m\), together with \(m\le M\), implies \(\gamma h\le1/4\), as needed for the local constant. The scalar recurrence theorem then gives \(e_n\le(1-3mh/(4\gamma))^ne_0+\gamma Mh\sqrt p/m\).



The actual numerical and reference position random variables provide a finite-cost Wasserstein coupling. The reference position law is the stated target \(\pi\), and position projection gives \(W_2\le\sqrt2e_n/\gamma\). Equal initial velocities identify \(e_0=\gamma\|X_0-Y_0\|_{L^2}\). Lean cancels friction and combines \(\sqrt2\sqrt p=\sqrt{2p}\), proving

\[W_2(\mathcal L(X_n),\pi)\le\sqrt2\left(1-\frac{3mh}{4\gamma}\right)^n\|X_0-Y_0\|_{L^2}+\frac{Mh\sqrt{2p}}m.\]

The initial term is the cost of the supplied coupling. Obtaining the infimum over all initial couplings, and discharging reference stationarity from the actual Brownian kinetic process, remain separate probabilistic requirements.



The proof is factored through a shared theorem with an explicit reference velocity radius \(C\ge0\). A uniform bound \(\|V_t^{\rm ref}\|_{L^2}\le C\) supplies \(C+e_k\) in the local defect and the terminal bias \(Mh\sqrt2C/m\). The standard-Gaussian statement is a thin specialization with the radius \(\sqrt p\) proved from its marginal law. This same shared theorem applies to particle temperature \(1/N\), using the separately proved thermal radius \(\sqrt{(d_\theta+Nd_x)/N}\); there is no duplicate exponential-integrator argument.



\noindent Lean declarations: \texttt{lp\_klmc\_reference\_wasserstein\_bound\_of\_velocity\_norm}, \texttt{lp\_klmc\_reference\_wasserstein\_bound}.

\end{document}
